\documentclass[12pt]{article}
\usepackage[utf8]{inputenc}
\usepackage[T1]{fontenc}
\usepackage[margin=1in]{geometry}
\usepackage{setspace}
\usepackage{xcolor}
\usepackage[colorlinks=true, urlcolor=blue!60!black, linkcolor=blue!60!black]{hyperref}
\setlength{\parindent}{0em}
\setlength{\parskip}{0.6em}
\renewcommand{\baselinestretch}{1.0}
\begin{document}
\begin{center}
{\LARGE \textit{McDonald v. City of Chicago}, 561 U.S. 742 (2010)\par}
\vspace{0.5em}
{\large Case Brief --- Two-Page Edition\par}
\end{center}

\section{Memory Jogger}
By 5--4, the Court held that the Fourteenth Amendment makes the Second Amendment right of armed self-defense binding on the States---four Justices through the Due Process Clause and Justice Thomas through the Privileges or Immunities Clause (\href{https://www.courtlistener.com/opinion/149702/mcdonald-v-city-of-chicago/}{561 U.S. 742}).

\section{Detailed Case Facts}
After \textit{Heller} recognized an individual right to keep a handgun at home against federal law, Chicago residents including Otis McDonald, a community activist threatened by drug dealers, challenged Chicago's and Oak Park's effective bans on handgun possession. The case turned on whether the Second Amendment binds the States.

\section{Procedural History}
The district court dismissed, and the Seventh Circuit affirmed, holding itself bound by \textit{Cruikshank}, \textit{Presser}, and \textit{Miller v. Texas} though it called their reasoning ``defunct.'' The Supreme Court reversed and remanded.

\section{Judicial Votes}
Reversed, 5--4. Alito wrote for Roberts, Scalia, Kennedy, and Thomas in part and for a plurality of four on Due Process incorporation; Scalia concurred; Thomas concurred in part and in the judgment; Stevens dissented, as did Breyer with Ginsburg and Sotomayor.

\section{Holding}
The Second Amendment right recognized in \textit{Heller} is incorporated against the States because it is fundamental to our scheme of ordered liberty and deeply rooted in the Nation's history and tradition, and it applies to the States as to the federal government.

\section{Analysis of Opinions}
Alito found self-defense ``the central component'' of the right and traced its roots from the English Bill of Rights through Reconstruction efforts to protect armed freedmen; the plurality declined to revive the Privileges or Immunities Clause or to dilute the right for States. Thomas called substantive due process a ``legal fiction'' and relied on citizens' privileges. Stevens denied that handgun ownership is implicit in ordered liberty, and Breyer stressed federalism and judicial incompetence on gun policy.

\section{Examples: Future Applications}
\textbf{Same-Side:} a ban on keeping handguns in public-housing apartments is unconstitutional. \textbf{Opposite-Side:} bans on possession by convicted felons remain valid under \textit{Heller}'s assurances. \textbf{Fence-Sitter:} a handgun ban for lawful permanent residents, given that Thomas's clause protects only citizens.

\section{Critique}
Originalists welcomed the result but faulted the refusal to overrule \textit{Slaughter-House}. Progressives argue it entrenched \textit{Heller}'s ``law-office history'' and hobbled cities, and scholars note a ``voting paradox'': majorities rejected each clause as the basis for incorporation.

\section{Key Quotations}
\begin{quote}
``In \textit{Heller}, we held that the Second Amendment protects the right to possess a handgun in the home for the purpose of self-defense. Unless considerations of \textit{stare decisis} counsel otherwise, a provision of the Bill of Rights that protects a right that is fundamental from an American perspective applies equally to the Federal Government and the States. See \textit{Duncan}, 391 U.S., at 149, and n.~14. We therefore hold that the Due Process Clause of the Fourteenth Amendment incorporates the Second Amendment right recognized in \textit{Heller}.''
\end{quote}
\noindent --- Justice Alito, Plurality Opinion, 561 U.S. at 791.

\end{document}
